% TOP_PROBLEM: {"id": 30002282, "title": "Dual-Graph Generation of Tautological Relations", "category_id": 5, "category": {"id": 5, "name": "algebraic_geometry", "display_name": "Algebraic Geometry", "description": "Geometric objects defined by polynomial equations.", "slug": "algebraic-geometry", "order_index": 5, "created_at": "2026-07-31T15:26:25.671Z"}, "status": "partially_solved"}
% FIRST_POSTED: null
% ENTRY_KIND: statement_only
% !TeX program = lualatex
\documentclass[11pt]{article}
\usepackage[margin=25mm]{geometry}
\usepackage{fontspec}
\setmainfont{Latin Modern Roman}
\usepackage{amsmath,amssymb,mathtools}
\usepackage{unicode-math}
\setmathfont{Latin Modern Math}
\usepackage{xurl,hyperref}
\hypersetup{colorlinks=true,urlcolor=blue}
\setcounter{secnumdepth}{0}
\setlength{\parindent}{0pt}
\setlength{\parskip}{5pt}
\setlength{\emergencystretch}{3em}
\newcommand{\R}{\mathbb R}
\newcommand{\N}{\mathbb N}
\newcommand{\Z}{\mathbb Z}
\newcommand{\Q}{\mathbb Q}
\newcommand{\C}{\mathbb C}
\newcommand{\D}{\mathbb D}
\newcommand{\F}{\mathbb F}
\newcommand{\sref}[2]{\hyperlink{source:#1:#2}{[S#2]}}
\newcommand{\cataloguescope}[1]{#1}
\begin{document}
% UnsolvedMath ID 30002282; source catalogue code OWR-12334-005
\setcounter{section}{30002281}
\section{Dual-Graph Generation of Tautological Relations}\label{30002282}
{\small UnsolvedMath ID 30002282 · Algebraic Geometry}
\subsection{Definitions and mathematical statement}

\paragraph{Shared notation.}$\mathbb N$, $\mathbb Z$, $\mathbb Q$, $\mathbb R$, and $\mathbb C$ denote the natural numbers, integers, rationals, reals, and complex numbers, respectively. All additional parameters and hypotheses are specified in the question.


\paragraph{Statement recorded in the supplied source.}
Let $\mathcal S_{g,n}$ be the strata algebra and let $\mathcal R\subset\mathcal S_{g,n}$ be the ideal generated by pullbacks and pushforwards of the dual-graph relations $R(g,n,r)$. Is $\mathcal R$ exactly the kernel of the natural surjection $\mathcal S_{g,n}\to R^*(\overline{\mathcal M}_{g,n})$? Do the analogous completeness statements hold after restriction to compact-type and rational-tail loci?
\par\smallskip\textit{Formulation references:} \sref{30002282}{1}, \sref{30002282}{2}.
\subsection{Short English statement}
Let $\mathcal S_{g,n}$ be the strata algebra and let $\mathcal R\subset\mathcal S_{g,n}$ be the ideal generated by pullbacks and pushforwards of the dual-graph relations $R(g,n,r)$. Is $\mathcal R$ exactly the kernel of the natural surjection $\mathcal S_{g,n}\to R^*(\overline{\mathcal M}_{g,n})$? Do the analogous completeness statements hold after restriction to compact-type and rational-tail loci?
\subsection{Sources}
\begin{itemize}
\item[\hypertarget{source:30002282:1}{S1}] ULAM.AI and the UnsolvedMath Contributors, UnsolvedMath: A Curated Collection of Open Mathematics Problems, record 30002282 (OWR-12334-005). Catalogue attribution under CC BY 4.0.
\url{https://huggingface.co/datasets/ulamai/UnsolvedMath}.
\item[\hypertarget{source:30002282:2}{S2}] Original problem formulation and mathematical context.
\url{https://doi.org/10.4171/owr/2013/06}.
\end{itemize}
\label{end:30002282}
\end{document}
